\section{Ô nhớ}
\label{sec:ch04-o-nho}

\index{ô nhớ!kiến trúc 1997}%
Bài đặt CEC vào giữa và dựng hai cổng nhân quanh nó. Kết quả gọi là
\tn{ô nhớ}{memory cell}. Cổng thứ nhất, \tn{cổng vào}{input gate}, quyết định
lúc nào input được ghi vào; cổng thứ hai, \tn{cổng ra}{output gate}, quyết định
lúc nào nội dung được đọc ra. Mỗi cổng là một unit sigmoid bình thường, có
trọng số riêng, huấn luyện bình thường. Cái không bình thường là cách dùng
output của chúng: nhân, chứ không cộng.

Viết theo ký hiệu của sách. Ba tiền kích hoạt, mỗi cái một khối trọng số:

\begin{align}
  \vect{a}^{i}_t &= \mat{W}_{xi}\vect{x}_t + \mat{W}_{hi}\vect{h}_{t-1} + \vect{b}_i,
  \label{eq:ch04-a-i} \\
  \vect{a}^{o}_t &= \mat{W}_{xo}\vect{x}_t + \mat{W}_{ho}\vect{h}_{t-1} + \vect{b}_o,
  \label{eq:ch04-a-o} \\
  \vect{a}^{c}_t &= \mat{W}_{xc}\vect{x}_t + \mat{W}_{hc}\vect{h}_{t-1} + \vect{b}_c.
  \label{eq:ch04-a-c}
\end{align}

Hai cổng, rồi trạng thái, rồi output:

\begin{align}
  \vect{i}_t &= \sigma(\vect{a}^{i}_t), \qquad
  \vect{o}_t = \sigma(\vect{a}^{o}_t),
  \label{eq:ch04-cong} \\
  \vect{c}_t &= \vect{c}_{t-1} + \vect{i}_t \odot g_{\mathrm{in}}(\vect{a}^{c}_t),
  \qquad \vect{c}_0 = \vect{0},
  \label{eq:ch04-cap-nhat-c} \\
  \vect{h}_t &= \vect{o}_t \odot g_{\mathrm{out}}(\vect{c}_t).
  \label{eq:ch04-h}
\end{align}

\index{trạng thái ô nhớ!hệ số bằng 1}%
Phương trình~\eqref{eq:ch04-cap-nhat-c} là chỗ cả chương xoay quanh. Hệ số của
\(\vect{c}_{t-1}\) là 1. Không phải một tham số khởi tạo gần 1, không phải một
đại lượng mạng tính ra: số 1, cố định, không học. Đó là
nghiệm~\eqref{eq:ch04-nghiem} của mục trước, đặt vào chỗ nó thuộc về.
\(\odot\) là nhân từng phần tử, và \tn{trạng thái ô nhớ}{cell state}
\(\vect{c}_t\) là đại lượng chạy dọc CEC.

Hình~\ref{fig:ch04-o-nho} vẽ một ô nhớ.

\begin{figure}[htbp]
  \centering
  % Ô nhớ LSTM 1997, vẽ lại hình 1 của bài báo theo ký hiệu của sách.
% Bản 1997: không có cổng quên, tự nối cố định 1.0.
% Mono-safe: chỉ dùng nét liền, nét đứt, mức xám và hình dạng node.
\begin{tikzpicture}[
  >= Stealth,
  ham/.style     = {draw, circle, minimum size = 7mm, inner sep = 0pt,
                    font = \footnotesize},
  nhan/.style    = {draw, circle, fill = black, minimum size = 2.2mm,
                    inner sep = 0pt},
  cong/.style    = {draw, circle, minimum size = 8mm, inner sep = 0pt,
                    font = \footnotesize},
  mui/.style     = {->, line width = 0.6pt, black!75},
  ghichu/.style  = {font = \footnotesize},
  nho/.style     = {font = \scriptsize},
]

  % --- đường chính: cell input -> g -> nhân -> CEC -> g_out -> nhân -> output
  \node[ham]  (g)    at (1.4, 1.6) {\(g_{\mathrm{in}}\)};
  \node[nhan] (nin)  at (2.9, 1.6) {};
  \node[ham]  (cec)  at (4.4, 1.6) {\(\vect{c}_t\)};
  \node[ham]  (gout) at (5.9, 1.6) {\(g_{\mathrm{out}}\)};
  \node[nhan] (nout) at (7.4, 1.6) {};

  % --- khung ô nhớ
  \begin{scope}[on background layer]
    \draw[black!55, line width = 0.5pt] (0.6, 0.75) rectangle (8.1, 2.75);
  \end{scope}
  \node[ghichu, anchor = north west, black!60] at (0.65, 2.72) {ô nhớ};

  % --- tự nối cố định 1.0: đây là CEC
  \draw[mui] (cec) to[out = 120, in = 60, looseness = 6]
    node[above, nho] {\(1.0\)} (cec);

  % --- luồng ngang
  \draw[mui] (0.0, 1.6) -- (g) node[midway, above, nho] {\(\vect{a}^{c}_t\)};
  \draw[mui] (g)    -- (nin);
  \draw[mui] (nin)  -- (cec);
  \draw[mui] (cec)  -- (gout);
  \draw[mui] (gout) -- (nout);
  \draw[mui] (nout) -- (8.7, 1.6) node[right, nho] {\(\vect{h}_t\)};

  % --- hai cổng, nằm dưới khung
  \node[cong] (cin)  at (2.9, -0.35) {\(\vect{i}_t\)};
  \node[cong] (cout) at (7.4, -0.35) {\(\vect{o}_t\)};
  \draw[mui] (cin)  -- (nin);
  \draw[mui] (cout) -- (nout);
  \draw[mui] (2.9, -1.35) -- (cin)
    node[midway, right, nho] {\(\vect{a}^{i}_t\)};
  \draw[mui] (7.4, -1.35) -- (cout)
    node[midway, right, nho] {\(\vect{a}^{o}_t\)};
  \node[nho, anchor = north] at (2.9, -1.4) {cổng vào};
  \node[nho, anchor = north] at (7.4, -1.4) {cổng ra};

  % --- chỗ trống mà bản 2000 sẽ lấp vào.
  % Đường dẫn đi thẳng xuống dưới hai nhãn cổng rồi mới có chữ: kéo ngang ở
  % cao độ của cổng thì chữ đâm vào chính node cổng ra.
  \draw[densely dashed, black!45, line width = 0.5pt] (4.4, 0.75) -- (4.4, -2.0);
  \node[nho, anchor = north, black!55, align = center] at (4.4, -2.05)
    {chỗ cổng quên của Gers 2000 sẽ nối vào;\\ bản 1997 để trống};

\end{tikzpicture}

  \caption{Một ô nhớ của bài 1997, vẽ theo hình 1 của bài. Đường ngang qua giữa
    là CEC: trạng thái \(\vect{c}_t\) mang tự nối cố định 1.0, vẽ thành vòng
    lặp phía trên nó, và đó là hệ số 1 trong~\eqref{eq:ch04-cap-nhat-c}. Hai
    cổng nằm dưới khung và tác động bằng phép nhân tại hai chấm đen: cổng vào
    chặn đường ghi, cổng ra chặn đường đọc. Nét đứt chỉ vào chỗ
    \tn{cổng quên}{forget gate} sẽ được lắp ba năm sau; bản 1997 không có gì ở
    đó.}
  \label{fig:ch04-o-nho}
\end{figure}

\subsection*{Hai hàm nén không phải tanh}

\index{hàm kích hoạt!g và h của bài 1997}%
Chỗ này là chỗ trí nhớ hay phản bội. Phụ lục A.1 của bài định nghĩa ba hàm, và
chỉ một trong ba là hàm bạn đoán:

\begin{align}
  \sigma(x) &= \frac{1}{1 + e^{-x}}, & &\text{khoảng giá trị } [0, 1],
  \label{eq:ch04-sigma} \\
  g_{\mathrm{out}}(x) &= \frac{2}{1 + e^{-x}} - 1, & &\text{khoảng giá trị } [-1, 1],
  \label{eq:ch04-g-out} \\
  g_{\mathrm{in}}(x) &= \frac{4}{1 + e^{-x}} - 2, & &\text{khoảng giá trị } [-2, 2].
  \label{eq:ch04-g-in}
\end{align}

Cả hai hàm nén là sigmoid co giãn, không phải tanh. Quan hệ đúng là
\(g_{\mathrm{out}}(x) = \tanh(x/2)\) và \(g_{\mathrm{in}}(x) = 2\tanh(x/2)\),
tức tanh với đối số bị chia đôi. Khoảng lệch giữa \(g_{\mathrm{in}}\) và
\(\tanh\) trên đoạn \([-6, 6]\) vượt 0.8 ở chỗ xa nhau nhất, nên đây không phải
sai số làm tròn mà là hai hàm khác nhau.

Một ghi chú về tên gọi. Bài viết hai hàm này là \(g\) và \(h\), nhưng \(h\) là
\tn{trạng thái ẩn}{hidden state} của sách, nên dùng nguyên tên của bài thì hai
ký hiệu đâm vào nhau ngay trong~\eqref{eq:ch04-h}. Sách viết
\(g_{\mathrm{in}}\) và \(g_{\mathrm{out}}\); phụ lục~\ref{app:ky-hieu} ghi lại
cả hai chiều.

\begin{minted}{python}
def g_in(x: torch.Tensor) -> torch.Tensor:
    """The paper's g, equation (5) of appendix A.1. Range [-2, 2]."""
    return 4.0 * torch.sigmoid(x) - 2.0


def g_out(x: torch.Tensor) -> torch.Tensor:
    """The paper's h, equation (4) of appendix A.1. Range [-1, 1]."""
    return 2.0 * torch.sigmoid(x) - 1.0
\end{minted}

Một bước của ô nhớ, từ repo companion, là bảy dòng:

\begin{minted}{python}
a_i = x @ self.w_xi.T + h_prev @ self.w_hi.T + self.b_i
a_o = x @ self.w_xo.T + h_prev @ self.w_ho.T + self.b_o
a_c = x @ self.w_xc.T + h_prev @ self.w_hc.T + self.b_c
i = torch.sigmoid(a_i)
o = torch.sigmoid(a_o)
c = c_prev + i * g_in(a_c)
h = o * g_out(c)
\end{minted}

\subsection*{Khối ô nhớ, và một chỗ sách đi chệch khỏi bài}

\index{khối ô nhớ!chia sẻ cổng}%
Bài không dừng ở một cell một cặp cổng. Nhiều cell dùng chung một cổng vào và
một cổng ra thì thành \tn{khối ô nhớ}{memory cell block}, và bài chọn cách này
vì lưu một thông tin phân tán vào một cell đơn lẻ không dễ. Khối cỡ 1 là một
cell thường; các thí nghiệm của bài dùng khối cỡ 1 và cỡ 2.

Còn một khác biệt lớn hơn, và nó lớn đến mức mục~\ref{sec:ch04-con-lai-gi} phải
quay lại đếm tham số theo cả hai cách. Trong bài, tầng ẩn nối \emph{đầy đủ}:
mỗi cổng nhận kết nối từ mọi ô nhớ \emph{và} từ mọi cổng khác. Thứ mà ngành
chốt lại về sau, và thứ \code{torch.nn.LSTM} cài đặt, là dạng tầng, nơi cả ba
khối trọng số đọc chung một \(\vect{h}_{t-1}\). Ba
phương trình~\eqref{eq:ch04-a-i} tới~\eqref{eq:ch04-a-c} ở trên là dạng tầng,
và repo companion cũng vậy.

Tôi chọn dạng tầng vì đó là thứ bạn sẽ gặp, không phải vì nó đúng với bài hơn.
Với mọi lập luận về đạo hàm trong chương này thì hai topology cho cùng kết
luận, vì cả hai đều không có đường nào chạy qua \(\vect{c}_{t-1}\) ngoài số
hạng đầu của~\eqref{eq:ch04-cap-nhat-c}. Chỗ chúng khác nhau là số tham số, và
đó là chỗ tôi sẽ nói lại.
